\part{Asynchronous Circuits}
\begin{nonumberframe}
    \vfill
    \centering
    \begin{beamercolorbox}[sep=8pt,center,shadow=false,rounded=true]{title}
    \usebeamerfont{title}Part II - ASYNCHRONOUS CIRCUITS\par%
    \usebeamerfont{subtitle} GUT meeting 13/01/2026
    \end{beamercolorbox}
    \vfill
\end{nonumberframe}

\section{Linear pipeline or digital Fordism}
\begin{frame}{Assembly Line}{\insertsection}
    % Figure Assembly line
\end{frame}

\begin{frame}{Linear Pipeline}{\insertsection}
    % Figure Combinational logic and registers
\end{frame}

\subsection{Concepts}
\begin{frame}{Latency \& Throughput}{\insertsection ~- \insertsubsection}
    % Definition Latency
    % Definition Throughput
\end{frame}

\begin{frame}{Parallelism \& Concurrency}{\insertsection ~- \insertsubsection}
    % Difference between both concepts
    % Introduction of a Fork and a Join
\end{frame}

\section{Clock-driven}
\begin{frame}{Clock Signal}{\insertsection}
    % Global Clock signal
    % Data propagation through global synchronization
\end{frame}
\subsection{Timing Constraints}
\begin{frame}{Linear Pipeline}{\insertsection ~- \insertsubsection}
    % Concept Critical Path
    % Worst case timing
\end{frame}

\begin{frame}{Setup}{\insertsection ~- \insertsubsection}
    % Definition of setup time and chronogram
\end{frame}

\begin{frame}{Hold}{\insertsection ~- \insertsubsection}
    % Definition of hold time and chronogram
\end{frame}

\subsection{Drawbacks}
\begin{frame}{Clock Skew}{\insertsection ~- \insertsubsection}
    % Problem with large clock tree
\end{frame}

\begin{frame}{Clock Domain Crossing}{\insertsection ~- \insertsubsection}
    % Synchronization of datas
    % GALS
\end{frame}

\begin{frame}{Robustness}{\insertsection ~- \insertsubsection}
    % Process variation, Voltage and Temperature
\end{frame}

\section{Event-driven}
\subsection{[TODO] Communication}
\begin{frame}{Local Synchronization}{\insertsection ~- \insertsubsection}
    % Pipeline with local synchronization
\end{frame}

\begin{frame}{Communication Channels}{\insertsection ~- \insertsubsection}
    % Sender, Receiver, Channel, Transaction
\end{frame}

\begin{frame}{Handshakes}{\insertsection ~- \insertsubsection}
    % Level / transition signaling
\end{frame}

\begin{frame}{Control Primitives}{\insertsection ~- \insertsubsection}
    % Uncond and Cond
\end{frame}

\subsection{Transactions at the gate level}
\begin{frame}{Delay Insensitive Circuits}{\insertsection ~- \insertsubsection}
    % Uncond and Cond
\end{frame}

\begin{frame}{State-Holding gates}{\insertsection ~- \insertsubsection}
    % Mueller gate / C-element
\end{frame}

\subsection{Formalism}
\begin{frame}{Petri Nets}{\insertsection ~- \insertsubsection}
    % Token, Place, Transition
\end{frame}

\begin{frame}{Signal Transition Graph}{\insertsection ~- \insertsubsection}
    % 4P Handshake and WCHB controller
\end{frame}

\section{Bundled-Data}
\begin{frame}{Transaction at the pipeline level}{\insertsection ~- \insertsubsection}
    % Control path, data path
    % matched delay
\end{frame}

\subsection{Timing Constraints}
\begin{frame}{Setup \& Hold}{\insertsection ~- \insertsubsection}
    % Relative Timing Constraints
    % Delay equations
\end{frame}

\begin{frame}{Local Clock Sets}{\insertsection ~- \insertsubsection}
    % Static Timing Analysis
    % During logical synthesis and PnR
\end{frame}

\section*{Summary}
\begin{frame}{\insertsection}

\end{frame}